\documentclass[aspectratio=169,12pt]{beamer}
\usepackage{amsmath}
\usepackage{amssymb}
\usepackage{array}
\usepackage[most]{tcolorbox}
\usepackage{graphicx}
\usepackage[sfdefault,lf]{FiraSans}
\renewcommand{\familydefault}{\sfdefault}
\setlength{\parindent}{0pt}
\everymath{\displaystyle}

\setbeamertemplate{navigation symbols}{}
\setbeamertemplate{footline}{}
\setbeamertemplate{headline}{}
\setbeamertemplate{blocks}[rounded][shadow=false]

\definecolor{mmpageWhite}{HTML}{FCFBF7}
\definecolor{mmtanSoft}{HTML}{F7F5EF}
\definecolor{mmgreenDeep}{HTML}{244B2C}
\definecolor{mmgreenLeaf}{HTML}{5D8A56}
\definecolor{mmgreenSoft}{HTML}{E4EFD9}
\definecolor{mmtext}{HTML}{163221}
\definecolor{mmwarn}{HTML}{8F2F36}
\definecolor{mmwarnSoft}{HTML}{F8E8EA}

\setbeamercolor{background canvas}{bg=mmpageWhite}
\setbeamercolor{normal text}{fg=mmtext,bg=mmpageWhite}
\setbeamercolor{itemize item}{fg=mmgreenDeep}
\setbeamercolor{itemize subitem}{fg=mmgreenDeep}

\newtcolorbox{MMCard}[2][]{enhanced,breakable=false,colback=#2,colframe=black,boxrule=1.5pt,arc=8pt,left=5pt,right=5pt,top=4pt,bottom=4pt,before skip=0pt,after skip=0pt,#1}
\newcommand{\KoruIcon}{\raisebox{-0.2em}{\includegraphics[height=1.05em]{../../../manamaths/SITE/header-logo.png}}}
\newcommand{\NotesTitle}[1]{{\colorbox{mmtanSoft}{\parbox{0.975\linewidth}{\vspace{0.14em}\hspace{0.25em}{\LARGE\bfseries\textcolor{mmgreenDeep}{#1}\hfill\KoruIcon}\vspace{0.14em}}}\par\vspace{0.18em}{\color{mmgreenLeaf}\rule{\linewidth}{2.0pt}}\par\vspace{0.12em}}}
\newcommand{\SectionTitle}[1]{{\large\bfseries\textcolor{mmgreenDeep}{#1}}\par\vspace{0.04em}}
\newcommand{\GreenDivider}{\par\vspace{0.01em}{\color{mmgreenLeaf}\rule{\linewidth}{2.4pt}}\par\vspace{0.03em}}
\newcommand{\KeyIdea}[1]{\begin{MMCard}[title={\textbf{Key idea}},colbacktitle=mmgreenSoft,coltitle=mmgreenDeep]{mmtanSoft}\small #1\end{MMCard}}
\newcommand{\WorkedExample}[2]{\begin{MMCard}[title={\textbf{#1}},colbacktitle=mmgreenSoft,coltitle=mmgreenDeep]{white}\small #2\end{MMCard}}
\newcommand{\CommonMistake}[1]{\begin{MMCard}[title={\textbf{Common mistake}},colbacktitle=mmwarnSoft,coltitle=mmwarn]{white}\small #1\end{MMCard}}
\newcommand{\TryThese}[1]{\begin{MMCard}[title={\textbf{Try these}},colbacktitle=mmgreenSoft,coltitle=mmgreenDeep]{mmtanSoft}\small #1\end{MMCard}}
\newcommand{\StepLine}[2]{\textbf{#1.} #2\par\vspace{0.03em}}

\usepackage{tikz}
\setbeamertemplate{background}{
 \begin{tikzpicture}[remember picture,overlay]
 \draw[black,line width=1.5pt,rounded corners=10pt]
 ([xshift=0.45cm,yshift=-0.45cm]current page.north west)
 rectangle
 ([xshift=-0.45cm,yshift=0.45cm]current page.south east);
 \end{tikzpicture}
}

\begin{document}

\begin{frame}[t,shrink=10]
\NotesTitle{Notes \& Steps}
\footnotesize
\begin{columns}[T,onlytextwidth]
\begin{column}{0.48\textwidth}
\KeyIdea{The power rule: if \(f(x)=x^n\) then \(f'(x)=nx^{n-1}\). Multiply by the exponent and reduce the exponent by one.}
\SectionTitle{Power rule table}
\[
\begin{array}{c|c}
f(x) & f'(x) \\ \hline
x^2 & 2x \\
x^3 & 3x^2 \\
x^4 & 4x^3 \\
x^{10} & 10x^9 \\
x^n & nx^{n-1}
\end{array}
\]
\SectionTitle{With a coefficient}
If \(f(x)=ax^n\) then \(f'(x)=anx^{n-1}\).
Example: \(f(x)=5x^4 \rightarrow f'(x)=20x^3\).
\end{column}
\begin{column}{0.48\textwidth}
\KeyIdea{Differentiation is linear: differentiate term by term and add the results. Constants differentiate to zero.}
\SectionTitle{Sum of terms}
\[
f(x)=3x^4+6x^3-7x^2-2x+3
\]
\[
f'(x)=12x^3+18x^2-14x-2
\]
\SectionTitle{Fractions as powers}
Write as powers of \(x\) first:
\[
\frac{1}{x}=x^{-1},\; \frac{1}{x^2}=x^{-2},\; \sqrt{x}=x^{\frac{1}{2}}
\]
Then apply the power rule.
\end{column}
\end{columns}
\GreenDivider
\CommonMistake{Forgetting to multiply by the coefficient too. For \(f(x)=3x^4\), the derivative is \(3 \times 4 \times x^3 = 12x^3\), not just \(4x^3\).}
\end{frame}

\begin{frame}[t,shrink=12]
\NotesTitle{Notes \& Steps}
\begin{columns}[T,onlytextwidth]
\begin{column}{0.48\textwidth}
\WorkedExample{Example 1: expand first}{Differentiate \(f(x)=(2x-1)(x+3)\).
\par Expand: \(2x^2+6x-x-3=2x^2+5x-3\).
\par Differentiate: \(f'(x)=4x+5\).}
\end{column}
\begin{column}{0.48\textwidth}
\WorkedExample{Example 2: fractions as powers}{Differentiate \(f(x)=3x^2+\dfrac{2}{x}+5\).
\par Rewrite: \(3x^2+2x^{-1}+5\).
\par Differentiate: \(f'(x)=6x-2x^{-2}=6x-\dfrac{2}{x^2}\).}
\end{column}
\end{columns}
\GreenDivider
\begin{columns}[T,onlytextwidth]
\begin{column}{0.48\textwidth}
\WorkedExample{Example 3: roots}{Differentiate \(f(x)=4\sqrt{x}-\dfrac{3}{x^2}\).
\par Rewrite: \(4x^{\frac{1}{2}}-3x^{-2}\).
\par Differentiate: \(f'(x)=2x^{-\frac{1}{2}}+6x^{-3}\).
\par So \(f'(x)=\dfrac{2}{\sqrt{x}}+\dfrac{6}{x^3}\).}
\end{column}
\begin{column}{0.48\textwidth}
\WorkedExample{Example 4: divide through first}{Differentiate \(f(x)=\dfrac{4x^3+7x^2-6x-5}{x^2}\).
\par Divide: \(4x+7-6x^{-1}-5x^{-2}\).
\par Differentiate: \(f'(x)=4+6x^{-2}+10x^{-3}\).
\par So \(f'(x)=4+\dfrac{6}{x^2}+\dfrac{10}{x^3}\).}
\end{column}
\end{columns}
\GreenDivider
\begin{columns}[b,onlytextwidth]
\begin{column}{0.48\textwidth}
\TryThese{\begin{enumerate}
  \item \(f(x)=6x^3-4x^2+2x-9\)
  \item \(f(x)=(x+4)(x-2)\)
  \item \(f(x)=5\sqrt{x}+\dfrac{1}{x^3}\)
\end{enumerate}}
\end{column}
\begin{column}{0.48\textwidth}
\CommonMistake{When rewriting fractions, remember each term is divided by the denominator. For \(\frac{x^3+2x}{x}\), write \(x^2+2\), not \(x^3+2\).}
\end{column}
\end{columns}
\end{frame}

\end{document}
